\documentclass[svgnames]{article}     % use "amsart" instead of "article" for AMSLaTeX format
%\geometry{landscape}                 % Activate for rotated page geometry

%\usepackage[parfill]{parskip}        % Activate to begin paragraphs with an empty line rather than an indent

\usepackage{graphicx}                 % Use pdf, png, jpg, or eps§ with pdflatex; use eps in DVI mode

%maths                                % TeX will automatically convert eps --> pdf in pdflatex
\usepackage{amssymb}
\usepackage{amsmath}
\usepackage{esint}
\usepackage[margin = 15mm]{geometry}

% Inverting Color of PDF
%\usepackage{xcolor}
%\pagecolor[rgb]{0.19,0.19,0.19}
%\color[rgb]{0.77,0.77,0.77}

%noindent
\setlength\parindent{0pt}

%pgfplots
\usepackage{pgfplots}

%images
\graphicspath{{/Users/devaldeliwala/downloads/}}                   % Activate to set a image directory

%tikz
\usepackage{pgfplots}
\pgfplotsset{compat=1.15}
\usepackage{comment}
\usetikzlibrary{arrows}
\usepackage[most]{tcolorbox}

%Figures
\usepackage{float}
\usepackage{caption}
\usepackage{lipsum}
\usepackage{subcaption}


\title{Detecting icy complex organic molecules and ions of Young Protostars}
\author{Deval Deliwala}
%\date{}                              % Activate to display a given date or no date

\begin{document}
\maketitle
%\section{}
%\subsection{}
\tableofcontents                     % Activate to display a table of contents

\section{Background Info}

Astronomical molecules provide insight into the physical conditions and
lifetimes of the sources they are in. They are associated with \textit{dense}
and \textit{cool} interstellar and circumstellar matter, which exist in either
the solid or gas phase of tiny dust particles. Non-terrestrial molecules are of
particular interest for what they tell us about the build-up of molecular
complexity throughout the universe. \\

When available in primitive planetary systems, COMs can potentially promote the
habitability of planets. 

\subsection{Utility} 

The utility of molecules as probes derives from both their spectra and
chemistry. 

\begin{itemize}
  \item[-] High-resolution rotational and vibrational spectra tell us about the
    density \& temperature of the gas and large-scale motions like
    collapse and rotation. 
  \item[-] Vibrational spectra of molecules in dust grains yield information on
    the polar/nonpolar nature of the mantles surrounding the particles.
\end{itemize}

\subsection{Complex Molecules} 

At present, around 150 different astronomical molecules have been detected. 50
or so are molecules of six or more atoms. These 50 molecules all contain at
least one carbon atom and are known as \textit{complex} molecules, or
\textit{complex organic molecules} (COMs). \\

COMs can be divided into two categories: 

\begin{itemize}
  \item[-] Unsaturated: organic species containing few hydrogen atoms 
    \item[-] Saturated: organic species rich in hydrogen, limited to chemical
      bonds involving only valence electrons. 
\end{itemize}

\textit{Most common terrestrial organic molecules tend to be saturated or
near-saturated.} \\

COMs have been extensively detected in the gas phase, thought to exist mostly
on icy grains . However, there have not been \textit{any} solid detections of
large-sized COMs, except for methanol ($CH_3OH$). 

\section{The Paper}

The goal of this paper is to ``comprehensively explore the COMs ice signature
in one low- and one high-mass protostar: NGC 1333 IRAS 2A and IRAS 23385+6053,
respectively."\\

It is only if organic material exists in ices (solid matter) as part of icy
dust grains that it would be effective in delivering water and organic
molecules to terrestrial planets. 

\subsection{Introduction}

The general consensus is that COMs form \textit{in the solid phase} and desorb
into a gas via thermal and non-thermal processes, particularly ion-molecule
collisions. \\

In the mid-IR spectral range, icy COMs have absorption features at $\sim 5.8\mu
m$, $6.8-7.0\mu m$, $7.2\mu m$, $7.4\mu m$, $\sim 8\mu m$, and $\sim 9.5-9.8\mu
m$. However, these bands overlap with features of other molecules. The bands at
$7.2\mu m$ and $7.4\mu m$ are consequently considered the \textit{main
fingerprint} for COMs. \\

This paper addresses the presence of simple molecules, COMs, and ions in
protostellar ices for NGC 1333 IRAS 2A and IRAS 23385+6053. 

\begin{itemize}
  \item[-] IRAS 23385+6053: High-mass $\sim 220 M_\odot; d = 4.9kpc$ star-forming region
    embedded in its natal molecular cloud with maser emissions of $H_2O$ and
    $CH_3 OH$. 
  \item[-] NGC 1333 IRAS 2A: Low-mass Class 0 Protostar, $d = 299\pm 3pc$ in
    a protobinary system. This source also contains a gaseous glycoaldehyde
    $HCOCH_2OH$, an important sugar molecule that participates in the formation
    of ribose, a \textit{component of RNA}. 
\end{itemize}

\subsection{Observation} 

\begin{figure}[htp]
  \centering
    \includegraphics[width = 7.5cm]{1.jpeg}
    \includegraphics[width = 10.5cm]{screenshot.png}
    \caption{NGCC 133 (Left) and IRAS 23885+6063 (Right)}
\end{figure}

\paragraph{IRAS 23385} 

\textit{Two-point dither mode} in a \textit{four-point mosaic} surrounding the
central protostellar cluster. \\

The observation revealed two mid-infrared continuum sources. They encompassed
the flux of both sources across the full wavelength range $5\mu m \rightarrow
25\mu m$. 


\paragraph{IRAS 2A} 

\textit{Two-point dither mode} for a \textit{single pointing} on the
protostellar binary. \\

Only continuum emission related to the primary component of the binary is
detected. They assume any contribution of the second star is negligible. \\

For both observations they optimized the aperture to capture as much source
flux without including too much noise. 

\textit{From JWST MIRI}: 

The Two-Point dither mode provides 2 exposures separated on the array and
avoids placing the peak of the PSF along the same row or column. This pattern
allows for simple background subtraction and minimizes the effects of
bad pixels, rows, and columns. \\

The data then went through a bunch of processing with the JWST calibration
pipeline to correct for fringes. 

\begin{figure}[H]
  \centering
    \includegraphics[width = \linewidth]{screenshot 2.png}
    \caption{Spectrum of both protostar sources. Strong gas-phase emission
      lines are masked in both spectra. The shaded cyan area indicates the
    ``COMs region" that will be studied. The lines labeled $H_{2}O$ bending
  and libration refer to the wavelength region where different vibrational
modes of $H_{2}O$ are at play.}
\end{figure}

\subsubsection{Background Subtraction for IRAS 23385+6053}

To accurately measure the photometry and morphological properties of
astronomical sources, one requires an accurate estimate of the background,
which can come from both the sky and the detector. \\

Usually the pipelines somehow work it out but the paper mentions they did not
work for IRAS 23385+6053. \\ 

They performed their own background subtraction by extracting a spectrum from
the observation off-source from the main-infrared continuum sources in the same
location of the sky where the background flux was the lowest. 

\subsection{Methodology} 

\begin{figure}[H]
  \centering
    \includegraphics[width = 6cm]{screenshot 3.png}
    \caption{Methodology used in the paper}
\end{figure}

\subsubsection{Continuum subtraction and silicate removal} 

\paragraph{Continuum} \mbox{} \\

The continuum spectra is resultant spectra of a star when viewed directly,
without any intervening matter. This results in a smooth spectra that matches
a black body. I believe the authors wish to remove the continuum spectra to
focus directly on the shaded cyan ``COMs" portion of the spectra. \\

Often a guided polynomial function is used to trace the continuum spectral
energy distributions (SEDs) for the protostars. They use a third-ordered
polynomial function to perform the fit, where there is little or no overlap
with the ice absorption bands of interest. \\

Once the SED is accurately modeled, they convert the spectra of both protostars
to an \textit{optical depth scale} using the equation below: 

\[
\tau_\lambda = -ln \left( \frac{F_\lambda^\text{source} }{F_\lambda^\text{cont}
} \right), 
\] \vspace{3px}

where $F_\lambda^\text{source}$ is the source spectrum and
$F_\lambda^\text{cont}$ is the continuum SED.  

\paragraph{Silica} \mbox{} \\ 

They used the `optool' code to generate optical depth
spectra of a silicate-carbon mixture, the typical chemical species of
interstellar grains with appropriate volume fractions used for protostar
models. They also account for amorphous carbon, dust grain sizes using some
power laws I have not thoroughly researched. \\

They find their optical depth ratio between the 9.8 $\mu m$ and 18  $\mu m$
features are about twice as high as the normal ratio found in the interstellar
medium. They only mention grain growth could be a reason and say it's beyond
the scope of this work. \\ 

By removing the silicate bans somehow they reveal the $H_{2}O$ libration bands
as well as other vibrational modes for ammonia and methanol. 

\begin{figure}[H]
  \centering
    \includegraphics[width = \linewidth]{screenshot 4.png}
    \caption{Continuum (a, b) and silicate subtraction (c,d) for IRAS 23385
      (left) and IRAS 2A (right). e and f show the silicate subtracted spectra
    of both stars with important features.} 
\end{figure}

\subsubsection{$SO_2$ gas emission subtraction} 

In the spectra of IRAS 2A, clear molecular emission is superimposed on the ice
absorption features between 7.25 $\mu$ m and 7.45 $\mu$ m. The emission was
found to be warm gas-phase $SO_2$. \\

They used a best-fit gas-phase emission line model of $SO_2$ to subtract the
$SO_2$ emission spectrum. This revealed another ice absorption feature around
7.4 $\mu$ m that was hidden previously. 

\subsubsection{Isolating ice features between 6.85 $\mu$ m and 8.6 $\mu$ m}


They use another third-order polynomial fit on the optical depth scale they
derived after removing the silicate. They then used `guiding-points' fitted by
another third-order polynomial function. They somehow also consider the $O-H$
vibrational mode that occurs in this wavelength range. 


\begin{figure}[H]
  \centering
    \includegraphics[width = \linewidth]{screenshot 5.png}
    \caption{Left: local continuum spectra 6.8-8.6 spectral range. Right:
    isolated optical depth spectra in the same range.}
\end{figure}

\subsubsection{Fits and lab data} 

They use a giant algorithm called `ENIIGMA' to search for the optimal
coefficients of the linear combinations of laboratory ice data. They thoroughly describe
how ENIIGMA works.

\begin{figure}[H]
  \centering
    \includegraphics[width = \linewidth]{screenshot 6.png}
    \caption{ENIIGMA fits for IRAS 23385 (top) and IRAS 2A (bottom).}
\end{figure}




\subsubsection{Degeneracy analysis} 

I am not entirely sure what they mean by degeneracy of the linear coefficients in this paper. It may just
be the energy degeneracy of different states of the molecules they find? I am
not sure. They implement a gaussian fit on each coefficient which allows them
to derive their confidence intervals. 

\subsubsection{Criteria for firm COM ice detections} 


\begin{itemize}
  \item[1.] Secure: multiple bands or bands of isotopologues are
detected in high-quality spectra; 
\item[2.] Likely: A single band is detected and the profile matches the
laboratory spectra; 
\item[3.] Possibly: Possibly: A single band is detected and there is no exact
match between the profile and the laboratory spectra.
\end{itemize}


\subsection{Results} 

\subsubsection{Spectral Decomposition \& Feature Analysis} 

The following COMs were fitted by ENIIGMA: 
\begin{itemize} 
  \item[1.] Complex Molecules
\begin{itemize}
  \item $CH_3CHO$ -- Robustly Detected
    \item $CH_3CH_2OH$ -- Robustly Detected
      \item $CH_3OHCHO$ -- Robustly Detected
        \item $CH_3COOH$
\end{itemize}
\item[2.] Simple Molecules
\begin{itemize}
  \item $CH_4$ 
    \item $SO_2$
    \item $HCOOH$
\end{itemize}
\end{itemize}

They also found a good match of $HCOOH^-$ and $OCH^-$ ions. \\

They continue describing which regions correspond to parts of the graph and why
these molecules likely emerged and where. Refer to Figure 6. \\

The recurrence of
each molecule between both protostars is shown below:  


\begin{figure}[H]
  \centering
    \includegraphics[width = 10cm]{screenshot 8.png}
    \caption{Recurrence of different chemical species for both protostars.}
\end{figure}

They afterwards test a lot more complex alcohols and other molecules. They go
deep into chemical analysis which I am having difficulty following. They
mention how if $H_2O$ is the major ice component, it means the COM features are
under the influence of the water ice polarity. They talk about generating an
`apolar ice matrix'. \\

They also definitively identify ions $HCOO^-$ and $OCN^-$ in both protostars. 

\subsubsection{Similarities with Comet 67P/G-C}

The main conclusion is that there is a good correlation between the abundances
of gas-phase COMs in high and low mass star-forming regions. However, COMs
abundance on the Comet are 10x larger. They mention this is because of ice
inheritance followed by `chemical alteration' toward later protostellar phases. 

\subsection{Conclusions}

\begin{itemize}
  \item[-] $CH_4$ ice is the main carrier of the 7.7 $\mu$ m band. $OCN^-$ and
    $SO_2$ contribute as well, which is the first documented time of an ion
    being attributed to this band. 
  \item[-] The 7.2 $\mu$ and 7.4 $\mu$m bands are mainly due to the format ion
    $HCOO^-$. This ion is formed via acid-base reactions, which are apparently
    common in interstellar ices. 
  \item[-] They confirm that $CH_3CH_2OH$ (ethanol) and $CH_3CHO$ (acetaldehyde)
    contribute to absorption lines and the 7.2 and 7.4 $\mu$ m, respectively.
    These molecules also contribute to other bands. 
  \item[-] They robustly detect $CH_3OCHO$ in both protostellar ices. They also
    `likely' find $CH_3COOH$ (acetic acid) in IRAS 2A. 
  \item[-] They show that the COMs found fit better in a polar environment,
    specifically when mixed in with water. 
  \item[-] The comparison with the comet strongly suggests that COMs in comets are
    significantly inherited from earlier protostellar phases. 
\end{itemize}



\end{document}

